% You can choose whether you prefer a single or double column appendix.
% Whatever you choose, you will need to stick to it throughout the appendix.
% \onecolumn

\chapter{First Appendix}
\label{sec:apx:first_appendix}

% Everything here is reference material the main text points at but does not need inline.
% The thesis must remain readable without it.

\section{Synthesis Algorithm Configurations}
\label{sec:apx:algorithms}
% Referenced from \ref{sec:method:data:generation:algorithms}. Reproduce, per algorithm
% (Orchestrate, Deepsyn, Syn4, C2RS): the exact abc command template, the parameters that
% differ from abc's defaults, and the alias definitions -- C2RS in particular expands via
% abc.rc into a long balance/resubstitute/rewrite chain that is worth showing in full.
% Source: AIG_DATASET_README.md and abc.rc. Pin the abc version here too; the labels depend
% on it and \ref{sec:discussion:limitations:reproducibility} flags that it is unrecorded.

\section{Reduction Method Parameters}
\label{sec:apx:params}
% One table: every reduction method, its configured parameters, and whether each was
% calibrated against measured retention or set to a plausible default. The distinction
% matters for \ref{sec:results:rq4:gap} -- an uncalibrated knob is a confound in a
% matched-compression comparison -- and putting it in one place keeps the methodology
% chapter from having to caveat each method individually.

\section{Baseline Hyperparameters}
\label{sec:apx:baselines}
% For SynthNet, HOGA and DeepGate4: which hyperparameters are the papers' published values
% and which had to be assumed or adapted, with the reason for each deviation. The two
% substantive ones (node-budget batching plus gradient accumulation) are argued in
% \ref{sec:method:architecture:baselines}; this is where the full table belongs.
% Source: src/train_baseline.py module docstring and src/baselines/*/regressor.py.

\section{Extended Results}
\label{sec:apx:results}
% Per-configuration tables too detailed for Chapter \ref{sec:results}: full metric sets per
% reduction method, per-tier error breakdowns, and the offline retention statistics per
% method and shard.